% !PRINTED 97
\weberpage{97}
%<<CONT
and for this, making use of the notation (6), we can set

\weq{8}{(\varrho, \eta) = \overset{n}{\Sigma}\, \varrho^{\sigma} r^n,}

wherein $n$ runs through all the numbers of the group $\mathfrak{N}$, and in $\sigma$ the
indices of $n$ are to be substituted each time for $\nu_1, \nu_2 \ldots \nu_\mu$.

Now let $r_1, r_2, \ldots r_\mu$ denote primitive roots of unity of the degrees
$q_1, q_2, \ldots q_\mu$, and let

\[ r = r_1 r_2, \ldots r_\mu \]

be set. Then the expression (8) can be decomposed into a product,
namely, if we set

\[ n \equiv n_h \ (\mathrm{mod}\ q_h) \]

into:

\weq{9}{(\varrho, \eta) = \Sigma\, \varrho^{\varepsilon_1 \nu_1} r_1^{n_1}\, \Sigma\, \varrho^{\varepsilon_2 \nu_2} r_2^{n_2} \ldots \Sigma\, r_\mu^{\varepsilon_\mu \nu_\mu} r_\mu^{n_\mu}.} % [sic: Weber prints r_\mu as the final first base; the parallel factors require \varrho.]

The factors of this product are now precisely the resolvents which we
investigated in \S.\,172 of the first volume, and we obtain, if we

\weq{10}{\overset{n_1}{\Sigma}\, \varrho^{\nu_h} r^{n_h} = (\varrho, \eta_h), \quad h = 1, 2, \ldots \mu} % [sic: Weber prints the upper bound n_1 and the base r^{n_h}; the general h factor requires n_h and r_h^{n_h}.]

set:

\weq{11}{(\varrho, \eta) = (\varrho^{\varepsilon_1}, \eta_1)(\varrho^{\varepsilon_2}, \eta_2) \ldots (\varrho^{\varepsilon_\mu}, \eta_\mu).}

If we now apply the formulae (5), (6), (25), (26) of the section just
cited, we obtain:

\weq{12}{(\varrho, \eta)(\varrho^{-1}, \eta) = m,}

\weq{13}{(\varrho, \eta)^3 = m(a + b\varrho), \quad (\varrho^2, \eta)^3 = m(a + b\varrho^2),}

wherein $a$ and $b$ are certain integers, which can be computed from the
formulae mentioned, but which turn out differently according to the choice
of the $\varepsilon$.

From (12) and (13) there then follows further:

\weq{14}{m = a^2 - ab + b^2.}

According to this we can set up the cubic equation whose roots are the
quantities $\eta, \eta', \eta''$. Let it read:

\weq{15}{\eta^3 - \alpha \eta^2 + \beta \eta - \gamma = 0,}

wherein $\alpha, \beta, \gamma$ are integers. As regards $\alpha$ first of
all, it is equal to the sum of all primitive $m$th
roots of unity, hence

\[ \alpha = \Sigma\, r_1\, \Sigma\, r_2 \ldots \Sigma\, r_\mu, \]

and this value is $= 0$ or $= (-1)^{\mu}$,
according as 9 occurs among the factors of $m$ or not. The two
other coefficients can be computed in just the same way as at the
%CONT>>
